\chapter{Hipótesis}
La Universidad Autónoma Chapingo es un lugar óptimo para construir una \texttt{Red de Sistemas de Captación de Agua de Lluvia} (RSCALL), que permitirá la sustentabilidad, ahorro y consumo de este recurso para la comunidad.
La creación de una red de sistemas es probable que sea más económica y fácil de operar que la construcción de cisternas localizadas en cada edificio;
Gracias al RSCALL, Chapingo podría reducir egresos económicos.
% Disponer de más agua, puede ser un incentivo para el desarrollo de sistemas hidropónicos y de agricultura vertical 

% Hablar sobre la calidad del agua de lluvia y su tratamiento. saber si sirve pa regar
\chapter{Justificación}
Cuánto se está gastando \$ y volúmenes de agua
Explicar lo de disponibilidad de agua en el futuro
Fomento de la sustentabilidad y uso consciente del agua
Innovación 
mantenimiento de las tuberías 
Cuánto podría costar si no lo hacen en diez años
Condiciones que tiene Texcoco
La necesidad de un modelo que pueda ajustarse en todos lados
\chapter{Revisión de Literatura}
Componentes básicos de un SCALL
Tratamiento de aguas de lluvia
Problemática del agua - INEGI
Clasificación de los sistemas de captación (humano, agrícola, ganadero, recarga de mantos acuíferos, agua de niebla),
normas y límites permisibles de la calidad de agua, análisis de la calidad del agua de lluvia contra la norma oficial del agua potable
redes de tuberías - hidráulica
uso para invernaderos







